\scanpage{16}
P.S. Here what are the modular varieties one gets from considerations of motive-theoretic galois theory.

\[
G\deleted{_{\mathbb{R}}}
= \text{reductive connected alg. group defined over }\mathbb{Q}
\]

\[
\begin{aligned}
U_{\mathbb{R}}={}&\text{dimension one twisted real Torus, with }U_{\mathbb{R}}(\mathbb{R})\\
&\text{identified with the group of complex nbs of}\\
&\text{module }1
\end{aligned}
\]

\[
\begin{aligned}
i_0:U_{\mathbb{R}}\longrightarrow G_{\mathbb{R}}\quad &\text{homomorphism of algebraic}\\
&\text{groups over }\mathbb{R},\ \text{s.th. }i_0(-1)\text{ is}\\
&\text{central, and s.th. \deleted{\unsure{t|a}}the centralisator}\\
&K_{\mathbb{R}}\text{ of }i_0(\sqrt{-1})\text{ is such that}\\
&K_{\mathbb{R}}(\mathbb{R})^{\unsure{0|\circ}}\text{ is \uline{maximal} \uline{compact} in }G_{\mathbb{R}}(\mathbb{R})^{0}.
\end{aligned}
\]

Moreover, for every $\mathbb{R}$-simple component of $G_{\mathbb{R}}$, the corresponding component of $i_0$ is not trivial.

\[
C_{\mathbb{R}}=\text{\uline{Centr}}_{G_{\mathbb{R}}}(i_0)\subset K_{\mathbb{R}}
\]

\[
\deleted{\unsure{\gamma|\mathscr{T}}}Q=G_{\mathbb{R}}/C_{\mathbb{R}}
\]

\[
\text{``Sigel-space''}=\unsure{\mathscr{J}|\mathcal{J}|\gamma}=G(\mathbb{R})^{0}/G(\mathbb{R})^{0}\cap C_{\mathbb{R}}(\mathbb{R})
\overset{\substack{\inserted{\text{(a connected component)}}\\ \deleted{\text{\unsure{open and closed}}}}}{\subset}Q(\mathbb{R})
\]

N.B. $\mathscr{J}$ has a natural complex structure (inherited from \inserted{one on} $Q(\mathbb{R})$ in fact) invariant by operations of $G(\mathbb{R})$. Every linear representation (over $\mathbb{Q}$) of $G$ in\deleted{to} a vector space $V$ (over $\mathbb{Q}$), together with a lattice $V_0$ of $V$, defines a `\uline{complex} \uline{analytic} family \inserted{$\unsure{\mathscr{V}|V}$} of Hodge structures\textsuperscript{\ensuremath{\delta}}\margin{$\delta$\;(\deleted{\unsure{perhaps even|perhaps}} \inserted{necessarily} globally polarizable -- to get a polarization of the family, one picks a ``form of polarization'' $\unsure{\varphi|\phi}$ \uline{on $V$} \deleted{\unsure{,|;}} invariant by $G$).} on $\mathscr{J}$, at least when one fixes a ``total degree'' $n$
% archive sheet number: 15
